% !TeX root = main page.tex
\documentclass[listof=totoc,a4paper,12pt,oneside,chapterprefix=true,sfdefaults=false]{scrbook}

% Template Designed by Dr. Alamgir Sardar (Assistant Professor), Department of Computer Science & Engineering, Sister Nivedita University, New Town, Kolkata, India. For any discrepancy email: alamgir.s@snuniv.ac.in

%======================Header Footer================================

% Load fancyhdr for custom headers/footers
\usepackage{fancyhdr}
\pagestyle{fancy}

% Clear default header and footer
\fancyhf{}

% ---------- HEADER ----------
% Display current chapter name in uppercase (like \headmark)
\fancyhead[RO,LE]{\nouppercase{\leftmark}}
\renewcommand{\headrulewidth}{0.4pt} % line under header

% ---------- FOOTER ----------
% Left (even pages): page number then a vertical rule then “Page”
\fancyfoot[LE]{\thepage\ \rule[-0.5ex]{1pt}{1.5ex}\ Page}

% Right (odd pages): “Page” then rule then page number
\fancyfoot[RO]{Page\ \rule[-0.5ex]{1pt}{1.5ex}\ \thepage}

% Center footer on all pages
\fancyfoot[L]{\textbf{\small Sister Nivedita University, WB}}

% ---------- PAGE STYLE ADJUSTMENTS ----------
% Make chapter-opening pages empty (no header/footer)
\fancypagestyle{plain}{
  \fancyhf{}
  \renewcommand{\headrulewidth}{0pt}
  \renewcommand{\footrulewidth}{0pt}
}

%==============================Header Footer=====================

\usepackage[margin=1in,head=27.2pt]{geometry}
% margin=1in → sets all four margins (left, right, top, bottom) to 1 inch.
% head=27.2pt → reserves 27.2 points of vertical space for the header. This adjusts how much space is allocated above the text block for running headers.

\usepackage[T1]{fontenc} % Controls font encoding (how characters are represented internally for typesetting)
% T1 is a European 8-bit encoding that supports accented characters (like é, ü, ñ).
% It also improves hyphenation for words with accented characters (compared to the default OT1 encoding).

% =============Text Formatting & Styles==========================
\usepackage[normalem]{ulem}    % underline, strikeout, emphasis (normalem avoids overriding \emph)
\usepackage{ragged2e}          % better text justification (e.g., \justifying)
\usepackage{textcomp}          % additional text symbols (e.g., °, €, etc.)
\usepackage{gensymb}           % generic symbols (e.g., degree, ohm)
\usepackage{csquotes}          % context-sensitive quotation marks
\usepackage{setspace}          % line spacing (single, 1.5, double)
% \usepackage{lineno}            % line numbering (disabled — unused, slow)
\usepackage{datetime}          % date and time formatting
\usepackage[none]{hyphenat}    % Disable hyphenation
% \usepackage{ulem}              % Underlining, strike-through, etc.
\usepackage{enumerate}         % Customizable enumerate environment
\usepackage{enumitem}          % Advanced control over lists (spacing, labels)
\usepackage{mfirstuc}          % Capitalize words / first letters


% ===================Mathematics Packages=================
\usepackage{amsmath}           % advanced math environments
\usepackage{amssymb}           % extra mathematical symbols
\usepackage{amsfonts}          % additional math fonts
\usepackage{mathrsfs}          % script fonts for math
\usepackage{mathtools}         % enhancements to amsmath
\usepackage[thinc]{esdiff}     % easy derivative notation
\usepackage{siunitx}           % consistent SI units and numbers
\usepackage{relsize}           %required for \mathlarger


% ==========Tables & Tabular Environments==============
\usepackage{tabularx}          % Tables with adjustable-width columns
\usepackage{longtable}         % multipage tables
\usepackage{multirow}          % multi-row cells in tables
\usepackage{multicol}          % multi-column formatting
\usepackage{mdwtab}            % IEEE-style tables
\usepackage{array}             % advanced array/tabular features
\usepackage{tabularray}        % Modern table package (like array + tabularx)
\UseTblrLibrary{booktabs}      % High-quality table rules for tabularray
\usepackage{booktabs}          % Professional-quality table rules
\usepackage{makecell}          % Line breaks & formatting inside table cells
% \usepackage{ragged2e}          % Better text justification in tables


% ===================Figures, Graphics & Diagrams=======================

\usepackage{graphicx}          % include graphics/images
\graphicspath{{./images/}}     % set path for images
% \usepackage[center]{caption}   % better caption formatting
\usepackage[justification=justified]{caption}   % better caption formatting
% \usepackage[justification=justified,singlelinecheck=false]{caption}
% \usepackage{bicaption}         % unused — bilingual captions not needed
\usepackage{subcaption}        % sub-figures/sub-tables with captions
% \usepackage{subfigure}         % Sub-figures (older, now replaced by subcaption)
\usepackage{float}             % improved float control (H option)
% \usepackage[inline]{asymptote} % disabled — runs external compiler, causes timeout
\usepackage{tikz}              % drawing diagrams
\usetikzlibrary{arrows,automata} % arrows, automata states
\usetikzlibrary{calc}          % coordinate calculations in TikZ
% \usepackage{overpic}           % unused
% \usepackage{epsfig,color}      % unused — color via xcolor if needed
\usepackage[figuresright]{rotating} % Rotate figures/tables
%\usepackage{psfig,epsfig}     % Obsolete graphics packages

% ==================BIBLIOGRAPHY / REFERENCES=====================
% \usepackage{natbib}     % For advanced citation styles (author–year, numeric, etc.)


% ==================DOCUMENT SPACING=============================

% \linespread{2}          % Double line spacing (useful for drafts/theses)

% =================Document Structure & Utilities=============================

\usepackage{pdfpages}          % include external PDF pages
% \usepackage{hyperref}          % hyperlinks and clickable references
\usepackage{hyperref,url}      % Hyperlinks for references, URLs, DOIs
\usepackage[section]{placeins} % prevents floats from crossing sections
\usepackage{eqparbox}          % equal-width boxes for text/math alignment
\usepackage[misc]{ifsym}       % additional symbols (envelopes, etc.)
\usepackage{mathptmx}          % Times Roman font for text + math


% ==================NOTES / ANNOTATIONS====================================
% \usepackage{todonotes}      % disabled — unused in final document

% ==================GLOSSARY===============================================
% \usepackage[nonumberlist,nogroupskip,xindy]{glossaries} % disabled — xindy is slow, not used

% ==================FOOTNOTES==============================================
\usepackage[bottom]{footmisc}  % Force footnotes to appear at bottom of page

% =================UTILITIES===============================================
\usepackage{etoolbox}          % Programming tools for LaTeX macros


% ==================ALGORITHMS==============================================
% \usepackage[linesnumbered,lined,commentsnumbered]{algorithm2e} 
% Pseudocode with numbered lines, comments, and formatting


\usepackage{algorithm,algpseudocode}
\usepackage{textgreek} % Allows Greek letters inside normal text without switching to math mode ($...$).

\renewcommand{\algorithmicrequire}{\textbf{Input: }}
\renewcommand{\algorithmicensure}{\textbf{Output: }}

% tikz already loaded above — libraries only:
\usetikzlibrary{positioning}
\usetikzlibrary{fit}
\usetikzlibrary{shapes.geometric}

% Following part to show colon after Algorithm #
% caption already loaded above — just add the setup:
\captionsetup[algorithm]{labelsep=colon}

%==========================================================================

\renewcommand{\labelitemi}{\tiny$ \color{black}\blacksquare$}  % blacksquare shaped itemized


\usepackage[backend=biber,style=numeric,sorting=none]{biblatex}
\addbibresource{references.bib}

\hyphenpenalty=1000000
\sloppy



% \renewcommand{\theadfont}{\normalsize}
\renewcommand\theadfont{\bfseries}  % For Bold Table head in center allignment write \thead{column name}
\renewcommand\theadgape{}
\newcolumntype{L}{>{\RaggedRight}X}
\newcolumntype{P}[1]{>{\RaggedRight\hspace{0pt}}p{#1}}

\newenvironment{calligraphic}{\usefont{T1}{pzc}{m}{it}}{}
\newenvironment{conditions}
  {\par\vspace{\abovedisplayskip}\noindent\begin{tabular}{>{$}l<{$} @{${}={}$} l}}
  {\end{tabular}\par\vspace{\belowdisplayskip}}
% Example


\hypersetup{
  colorlinks   = true, %Colours links instead of ugly boxes
  urlcolor     = blue, %Colour for external hyperlinks
  linkcolor    = blue, %Colour of internal links
  citecolor   = red %Colour of citations
}


% Define a new font command
\newcommand{\setchapterfont}{\fontsize{12}{14}\selectfont\fontfamily{ptm}\selectfont}
% Patch the \mainmatter command to change font settings
\preto{\mainmatter}{\setchapterfont}
\newcommand{\acfull}[1]{\acl{#1} (\acs{#1})}

% handy shortcut macro:for table in introduction
\newcommand\mytab[1]{%
   \begin{tabular}[t]{@{}l@{}} #1 \end{tabular}}


% ----------------------Paragraph Spacing-----------------------------

\setlength{\parskip}{6pt}   % space BETWEEN paragraphs
\setlength{\parindent}{0pt} % optional: remove paragraph indentation
\usepackage{parskip}


% ----------------------Remove space before chapter heading--------------

\makeatletter
\renewcommand{\@makechapterhead}[1]{%
    \vspace*{-20pt}% reduce space before chapter title
    {\parindent \z@ 
     \normalfont
     \huge\bfseries
     \ifnum \c@chapter > 0
       \chaptername\ \thechapter % prints "Chapter 1"
       \par\nobreak
       \vskip 10pt
     \fi
     #1\par}%
    \vspace{20pt}%
}
\makeatother
% ---------------------------------------------------------------------



\input{bibtext/class file}

\begin{document}

\frontmatter




\begin{titlepage}
% !TeX root = ../main page.tex


\begin{center}

    \vspace{5mm}

    \projectTitle{Development of a Multi-Modal Surveillance System Integrating Aerial Imaging and Signal-Based Detection for Enhanced Situational Awareness}

    \programName{Bachelor of Technology (Computer Science and Engineering)}



%---------Provide Student Name, Registration Number, EmailID here-------------
% If you have less group members then delete remaining rows from the following table

\begin{table}[H]
\centering
\renewcommand{\arraystretch}{1.4}
    \begin{tabular}{lcl}
    \textbf{Student Name} & \textbf{Reg. Number} & \textbf{Email ID}  \\ %\hline

    \textbf{TAMANNA DAS} & \textbf{240013403728} & \textbf{tamannadas052@gmail.com} \\

    \textbf{MITUDRU DUTTA} & \textbf{230014970503} & \textbf{mitudrudutta72@gmail.com} \\

    \textbf{KAUSHIK SAMADDER} & \textbf{230014989236} & \textbf{kaushiksamadder2003@gmail.com} \\

    \textbf{SIDHARTH PANDIT} & \textbf{230014342653} & \textbf{sidharthpandit063@gmail.com} \\

    \textbf{DEBANSHU PAUL} & \textbf{230015088078} & \textbf{pauldebanshu73@gmail.com} \\
    \end{tabular}
\end{table}


 \vspace{2cm}

     %---------------------Change date below------------------------
    \submissionDate{April 30, 2026}

    \vspace{1cm}

    % \supervisorName{Name of the Supervisor}


    % ----------------Change supervisor Designation if need----------------
    % \supervisorDesignation{Assistant Professor}




    \universityDetailsTitlePage
    \vspace{-0.6cm}


    \begin{center}
    \vspace{0.2cm}
        {\Large April, 2026 \par}
    \end{center}


\end{center}

\end{titlepage}

\setcounter{page}{1}



\setcounter{tocdepth}{1}
\tableofcontents
\listoftables
\listoffigures


\justifying
\onehalfspacing
% !TeX root = ../main page.tex
\chapter{Abstract}
\onehalfspacing

Cameras still dominate surveillance, yet pictures alone fail where visibility collapses---night patrols, smoke after a fire, fog near coastlines, or targets half-hidden behind trees and structures. Radar behaves differently: it keeps responding when optics cannot. The trade-off is meaning---radar rarely tells you \emph{who} or \emph{what} moved without heavy interpretation. Our goal is to bridge that gap without enterprise-grade budgets.

We develop a compact aerial prototype built around an F450 frame, Raspberry~Pi~4, standard Pi camera, LD2410 24~GHz radar, and GPS. On the drone we run YOLOv8-pose for people and posture cues, ByteTrack to hold identities across drops in visibility, and a Doppler--CFAR--MTI chain on radar samples so motion is not guessed from pixels alone. A small Multi-Layer Perceptron (MLP) merges radar and vision features into one threat score per track---simple enough to ship on Pi-class silicon.

Alerts go beyond bounding boxes: lightweight rules flag loitering, intrusion-like motion, crowded clustering, and possible falls; a Kalman filter rolls tracks forward a few seconds for warning lead time. Streams leave the aircraft over WebSocket/MQTT into a backend that feeds a command-centre-style dashboard and an AR handset app for whoever is on the ground. Stored events can be queried through an LLM-assisted interface so operators are not forced to write SQL mid-incident.

We plan to judge the stack honestly---vision-only vs.\ radar-only vs.\ fused paths---using precision, recall, $F_1$, latency, and robustness checks on data we collect ourselves. Success means a documented, inexpensive rig that survives bad weather better than cameras alone while staying reproducible for labs like ours.

\textbf{Keywords:} Multi-modal surveillance, sensor fusion, UAV, mmWave radar, edge AI.


\setlength{\parskip}{6pt}   % space between paragraphs

\mainmatter

\chapter{Introduction}


%----------Focus on the following points during discusion------------

% Explain the context of the problem
% State why the topic is important
% Mention current trends or scenario
% Give a brief overview of the domain









\chapter{Problem Statement}


%----------Focus on the following points during discusion------------

% Summarize existing research work
% Highlight key findings
% Identify limitations/gaps
% Show how your work is different or needed

\chapter{Significance of the Study}

%----------Focus on the following points during discusion------------

% discuss about Who benefits  
% how Academic / industrial / societal impact

\chapter{Literature Review}
\label{Literature_Review}

This chapter reviews representative research relevant to the proposed multi-modal aerial surveillance system. The discussion follows the style of paper-wise analysis and highlights how each prior work contributes to the problem space. The review is grouped into four themes: UAV vision-based detection, radar-based sensing, Wi-Fi CSI and passive RF sensing, and multi-modal fusion with edge artificial intelligence. The chapter concludes with a synthesis that directly motivates the proposed system architecture.

\section{UAV-Based Aerial Imaging and Detection}

Redmon et al.~\cite{redmon2016yolo} introduced the YOLO detection framework and showed that unified single-stage object detection can achieve real-time speed while maintaining strong accuracy. This work established the baseline for many practical surveillance systems where low latency is critical.

Jocher et al.~\cite{jocher2023yolov8} advanced this line with YOLOv8 and demonstrated improved training stability, faster inference variants, and broader task support including pose estimation. The pose branch is especially useful for behaviour-level surveillance because it captures posture cues rather than only bounding boxes.

Wojke et al.~\cite{wojke2017deepsort} proposed DeepSORT, combining Kalman filtering and appearance embeddings for identity-preserving multi-object tracking. Their method significantly reduced identity switches compared to motion-only association.

Zhang et al.~\cite{zhang2022bytetrack} introduced ByteTrack and showed that associating both high- and low-confidence detections improves trajectory continuity under partial occlusion. This is directly relevant to aerial scenes where targets are often intermittently visible because of viewpoint changes and clutter.

Although these works provide strong visual detection and tracking foundations, camera-only pipelines remain sensitive to adverse weather, illumination changes, smoke, and foliage occlusion. They also do not directly provide radar-like range-velocity cues, motivating the inclusion of complementary sensing.

\section{Radar-Based Detection}

Skolnik~\cite{skolnik2008radar} and Richards~\cite{richards2014principles} provide the theoretical basis for practical radar perception, including Doppler processing, Moving Target Indication (MTI), and Constant False Alarm Rate (CFAR) detection. These techniques remain central for extracting motion from cluttered environments.

Patole et al.~\cite{patole2017automotive} surveyed automotive radar pipelines and highlighted radar's robustness in poor visibility, especially when optical sensors degrade. Their review also emphasised the challenge of semantic classification using radar alone.

Santra and Hazra~\cite{santra2020short} demonstrated short-range mmWave radar for human motion and gesture understanding using micro-Doppler patterns. Their results confirmed that RF sensing can preserve detection capability in challenging visual conditions.

In practical low-cost systems, modules such as LD2410 make radar sensing accessible; however, radar-only pipelines still face limited spatial detail and weaker semantic interpretation than vision systems. This motivates feature-level fusion rather than single-modality operation.

\section{Wi-Fi Channel State Information and Passive RF Sensing}

Kotaru et al.~\cite{kotaru2015spotfi} proposed SpotFi and demonstrated device-free localization from Wi-Fi CSI using fine-grained channel information. Their results established CSI as a viable passive sensing signal beyond communication.

Zhang et al.~\cite{zhang2019farsense} presented FarSense and showed how CSI phase/amplitude fluctuations can be used for respiration estimation. This work is important because it demonstrates subtle physiological sensing without wearable devices.

Zheng et al.~\cite{zheng2019zero} introduced Widar 3.0 for cross-domain gesture recognition and improved generalisation across environments and users. Their study highlighted how domain adaptation is critical for robust RF sensing.

Recent open-source implementations such as RuView~\cite{ruview} show practical real-time CSI data pipelines, tracker integration, and edge-to-dashboard streaming patterns. Even when CSI is not used as the primary modality in outdoor UAV operation, these architectural and algorithmic patterns are valuable references for robust RF-aware system design.

Most CSI literature remains strongest in indoor/static deployments and is sensitive to multipath noise, which limits direct transfer to airborne outdoor surveillance. Therefore, CSI is treated in this work as conceptual and architectural guidance rather than the primary sensing channel.

\section{Multi-Modal Sensor Fusion and Edge Artificial Intelligence}

Khaleghi et al.~\cite{khaleghi2013multisensor} categorised multi-sensor fusion into early, feature-level, and decision-level strategies. Their taxonomy is used in this proposal to justify selecting feature-level fusion for balancing interpretability, computational cost, and performance.

Chadwick et al.~\cite{chadwick2019distantvehicle} demonstrated that combining radar and camera cues improves distant-object detection reliability over camera-only baselines. Their findings support the argument that fusion is most beneficial when one modality weakens.

Nabati et al.~\cite{nabati2021centerfusion} proposed CenterFusion and showed that radar-camera feature integration improves object detection quality and robustness in adverse scenarios. This work motivates the use of lightweight learned fusion in this proposal.

On the deployment side, open-source systems such as PLFM\_RADAR~\cite{plfm_radar} demonstrate accessible signal-processing workflows, while modern embedded hardware (for example Raspberry~Pi-class devices) enables practical edge inference for low-cost prototypes.

Despite progress, most published multi-modal systems either depend on expensive sensor stacks or focus on ground-vehicle contexts rather than UAV-centric field deployment with unified operator interfaces.

\section{Synthesis}

The reviewed works collectively indicate four clear conclusions:  
(i) vision models offer strong semantic understanding but are vulnerable in degraded visibility;  
(ii) radar sensing is robust under poor visual conditions but weak in semantic interpretation;  
(iii) CSI research provides useful RF sensing and pipeline ideas, but mostly in indoor settings; and  
(iv) fusion improves reliability, yet low-cost UAV-native end-to-end implementations remain limited.

Therefore, the proposed project positions itself as a practical integration effort: combining UAV-based imaging, low-cost mmWave radar, edge AI fusion, behaviour analysis, and real-time dashboard/mobile interfaces into one reproducible surveillance framework suitable for resource-constrained academic and operational contexts.

\chapter{Proposed Methodology}

\textit{\textcolor{red}{Write a brief idea about your proposed methodology then It is mandatory to Insert a clear Block Diagram/ system architecture diagram (which one is applicable) in this chapter. A sample example is shown below:}}


%----------Focus on the following points during discusion------------

% Describe the proposed approach or method
% Explain:
%     Tools/techniques used
%     Steps involved
% Include a diagram (mandatory):
%     Block Diagram / System Architecture / Workflow

\chapter{Research Gap Identification}
\label{Research_Gap_Identification}
    
%----------Focus on the following points during discusion------------

% Clearly define the problem you are addressing
% Identify specific gaps in current systems/research
% Ensure the gap is:
%     Precise
%     Researchable

\chapter{Objectives}
\label{Objectives}


%----------Focus on the following points during discusion------------

% Main objective (1–2 lines)
% Specific objectives (bullet points), such as:
% To analyze…
% To design…
% To implement…
% To evaluate…






\chapter{Work Plan}

%----------Focus on the following points during discusion------------

% Divide work into phases or semesters
% Example:

% Project-I

% Literature Review → Month 1
% Design & Implementation → Month 2–3
% Testing → Month 4–5

% Project-II

% Optimization → Month 1–2
% Final Testing → Month 3
% Documentation → Month 4

\backmatter

% \bibliographystyle{IEEEtran}
\printbibliography
\end{document} 